\documentclass[12pt, a4paper]{article}

\usepackage[utf8]{inputenc}
\usepackage[T1]{fontenc}
\usepackage[portuguese]{babel}
\usepackage{geometry}
\usepackage{graphicx}
\usepackage{hyperref}
\usepackage{listings}
\usepackage{xcolor}
\usepackage{booktabs}
\usepackage{array}
\usepackage{longtable}
\usepackage{float}
\usepackage{titlesec}
\usepackage{parskip}
\usepackage{fancyhdr}
\usepackage{enumitem}
\usepackage{mdframed}

\geometry{margin=2.5cm}

% Cores
\definecolor{pgblue}{RGB}{51, 103, 145}
\definecolor{codebg}{RGB}{245, 245, 245}
\definecolor{codeframe}{RGB}{200, 200, 200}
\definecolor{keywordcolor}{RGB}{0, 0, 180}
\definecolor{commentcolor}{RGB}{80, 120, 80}
\definecolor{stringcolor}{RGB}{160, 40, 40}
\definecolor{noteblue}{RGB}{230, 240, 255}
\definecolor{noteframe}{RGB}{100, 140, 200}
\definecolor{warnbg}{RGB}{255, 248, 230}
\definecolor{warnframe}{RGB}{210, 160, 50}

% Estilo SQL
\lstdefinelanguage{SQL}{
  keywords={CREATE, TABLE, PRIMARY, KEY, FOREIGN, REFERENCES, NOT, NULL,
            UNIQUE, INDEX, ON, DELETE, CASCADE, SET, DEFAULT, CHECK,
            INSERT, INTO, VALUES, SELECT, FROM, WHERE, AND, OR, IN,
            VARCHAR, INT, INTEGER, FLOAT, TEXT, TIMESTAMP, DATETIME,
            BOOLEAN, SERIAL, TYPE, ENUM, AS, CONSTRAINT, ALTER, ADD},
  keywordstyle=\color{keywordcolor}\bfseries,
  comment=[l]{--},
  commentstyle=\color{commentcolor}\itshape,
  stringstyle=\color{stringcolor},
  morestring=[b]',
  sensitive=false
}

\lstset{
  language=SQL,
  backgroundcolor=\color{codebg},
  frame=single,
  rulecolor=\color{codeframe},
  basicstyle=\ttfamily\small,
  breaklines=true,
  showstringspaces=false,
  tabsize=2,
  captionpos=b,
  numbers=left,
  numberstyle=\tiny\color{gray},
  numbersep=8pt,
  xleftmargin=15pt
}

% Caixas de Nota e Aviso
\newmdenv[
  backgroundcolor=noteblue,
  linecolor=noteframe,
  linewidth=1pt,
  innerleftmargin=10pt,
  innerrightmargin=10pt,
  innertopmargin=8pt,
  innerbottommargin=8pt,
  skipabove=10pt,
  skipbelow=10pt
]{notebox}

\newmdenv[
  backgroundcolor=warnbg,
  linecolor=warnframe,
  linewidth=1pt,
  innerleftmargin=10pt,
  innerrightmargin=10pt,
  innertopmargin=8pt,
  innerbottommargin=8pt,
  skipabove=10pt,
  skipbelow=10pt
]{warnbox}

% Cabeçalho / Rodapé
\pagestyle{fancy}
\fancyhf{}
\rhead{\textcolor{gray}{\small IDAP -- Design de Base de Dados}}
\lhead{\textcolor{gray}{\small ISEL -- Projeto e Seminário 2025/26}}
\rfoot{\textcolor{gray}{\small \thepage}}
\renewcommand{\headrulewidth}{0.4pt}

% Formatação de Títulos
\titleformat{\section}{\large\bfseries\color{pgblue}}{\thesection}{1em}{}[\titlerule]
\titleformat{\subsection}{\normalsize\bfseries}{\thesubsection}{1em}{}

% Configuração Hyperref
\hypersetup{
  colorlinks=true,
  linkcolor=pgblue,
  urlcolor=pgblue,
  citecolor=pgblue,
  pdftitle={IDAP - Documentação de Design de Base de Dados},
  pdfauthor={Afonso Santos \& Rodrigo Meneses}
}

\begin{document}

% Página de Rosto
\begin{titlepage}
  \centering
  \vspace*{2cm}
  {\Huge\bfseries IDAP\par}
  \vspace{0.4cm}
  {\Large Insurance Damage Assessment Platform\par}
  \vspace{0.3cm}
  \rule{0.6\linewidth}{0.5pt}
  \vspace{0.4cm}

  {\large\bfseries Documentação de Design de Base de Dados\par}
  \vspace{0.2cm}
  {\normalsize Modelo ER, Considerações de Design e Implementação em PostgreSQL\par}

  \vspace{2cm}

  \begin{tabular}{rl}
    \textbf{Alunos:}     & Afonso Santos (n.º 50484) \\
                         & Rodrigo Meneses (n.º 50542) \[0.5em]
    \textbf{Orientadores:}& Artur Ferreira \\
                         & Pedro Matutino \[0.5em]
    \textbf{Curso:}      & Licenciatura em Engenharia Informática e de Computadores \\
    \textbf{UC:}    & Projeto e Seminário -- Semestre de Verão 2025/26 \\
    \textbf{Data:}       & Março de 2026 \\
  \end{tabular}

  \vfill
  {\small Instituto Superior de Engenharia de Lisboa (ISEL)}
\end{titlepage}

\tableofcontents
\newpage

% -------------------------------------------------
\section{Modelo Entidade--Relacionamento}

\subsection{Visão Geral}

O modelo ER está organizado em torno do conceito central de \texttt{AccidentCase} (Processo de Sinistro), que agrupa todas as entidades envolvidas numa única peritagem de seguros. A partir de um processo, o modelo ramifica-se em quatro áreas principais: utilizadores e autenticação, veículos e danos, evidências e medições, e análise e relatórios.

\begin{notebox}
\textbf{Nota:} O diagrama detalhado do modelo Entidade-Relacionamento encontra-se disponível para consulta na última página deste documento (Anexo A). O diagrama ilustra a entidade central \texttt{AccidentCase} ligada às entidades de utilizadores, veículos, evidências e análise, refletindo as relações descritas na secção seguinte.
\end{notebox}

\subsection{Descrição das Entidades}

\begin{notebox}
\textbf{Nota:} É possível que vão existir mais tabelas derivadas de \textbf{Evidence} mas decidimos começar apenas pela análise de imagens.
\end{notebox}

\begin{longtable}{>{\bfseries}p{4cm} p{10cm}}
\toprule
Entidade & Descrição \\
\midrule
\endhead
Role & Tabela de consulta que define os perfis de utilizador: Supervisor, Gestor, Triador, Admin do sistema. \[0.4em]
User & Representa qualquer pessoa que se autentique na plataforma. Armazena credenciais (apenas hash da palavra-passe) e uma chave estrangeira para \texttt{Role}. \[0.4em]
TOKENS & Armazena tokens de autenticação ativos. Suporta gestão de sessões e revogação de tokens. \[0.4em]
AccidentCase & A entidade central. Representa um processo de sinistro individual. \[0.4em]
WeatherConditions & Regista o contexto meteorológico no momento do acidente (temperatura, visibilidade, tipo de precipitação). \[0.4em]
AccidentScene & Captura o contexto geométrico e ambiental do local: coordenadas GPS, inclinação do terreno, gradiente da estrada e tipo de via. \[0.4em]
Vehicle & Um veículo envolvido num processo. Identificado pela marca, modelo, ano e matrícula. \[0.4em]
Damage & Um dano específico observado num veículo. Regista a zona de contacto, tipo de deformação, direção do impacto e altura medida ao solo. \[0.4em]
Evidence & Qualquer prova anexada ao processo (foto, vídeo, documento). \[0.4em]
ImageEvidence & Extensão de \texttt{Evidence} para ficheiros de imagem. Armazena o caminho do ficheiro, dimensões e metadados (ex: dados EXIF). \[0.4em]
Analysis & Uma sessão de análise formal realizada por um perito sobre um processo. \[0.4em]
AnalysisImage & Tabela de junção que liga uma análise às imagens específicas utilizadas como prova. \[0.4em]
Measurement & Armazena as coordenadas de pixéis de um objeto de referência e de uma área de dano para calcular a altura real. \[0.4em]
DamageComparison & Regista a comparação formal entre dois danos de veículos diferentes. \[0.4em]
AnalysisConclusion & A conclusão estruturada de uma análise: resultado de compatibilidade e fundamentação do perito. \[0.4em]
Report & O documento final gerado para um processo. \\
\bottomrule
\caption{Descrição das entidades}
\end{longtable}

\subsection{Relações e Cardinalidades}

A tabela seguinte resume as principais relações entre entidades e as respetivas cardinalidades.

\begin{longtable}{>{\bfseries}p{3.5cm} p{1.8cm} >{\bfseries}p{3.5cm} p{5.5cm}}
\toprule
Entidade A & Cardinalidade & Entidade B & Descrição \\
\midrule
\endhead
Role & 1 : N & User & Um perfil pode ser atribuído a vários utilizadores; cada utilizador tem exatamente um perfil. \[0.4em]
User & 1 : N & AccidentCase & Um utilizador pode gerir vários processos de sinistro; cada processo é gerido por um utilizador. \[0.4em]
User & 1 : N & TOKENS & Um utilizador pode ter vários tokens de sessão ativos; cada token pertence a um utilizador. \[0.4em]
AccidentCase & 1 : 1 & WeatherConditions & Cada processo tem exatamente um registo de condições meteorológicas. \[0.4em]
AccidentCase & 1 : 1 & AccidentScene & Cada processo tem exatamente um registo do local do acidente. \[0.4em]
AccidentCase & 1 : N & Vehicle & Um processo pode envolver vários veículos; cada veículo pertence a um processo. \[0.4em]
Vehicle & 1 : N & Damage & Um veículo pode ter vários danos registados; cada dano está associado a um único veículo. \[0.4em]
AccidentCase & 1 : N & Evidence & Um processo pode ter várias evidências anexadas; cada evidência pertence a um único processo. \[0.4em]
Evidence & 1 : 1 & ImageEvidence & Uma evidência do tipo imagem tem exatamente uma entrada em \texttt{ImageEvidence} (herança de tabela de classe). \[0.4em]
AccidentCase & 1 : N & Analysis & Um processo pode ser alvo de várias sessões de análise; cada análise pertence a um processo. \[0.4em]
Analysis & N : M & ImageEvidence & Uma análise pode utilizar várias imagens; uma imagem pode ser usada em várias análises (via \texttt{AnalysisImage}). \[0.4em]
Analysis & 1 : N & Measurement & Uma análise pode conter várias medições; cada medição pertence a uma análise. \[0.4em]
Analysis & 1 : N & DamageComparison & Uma análise pode incluir várias comparações de danos. \[0.4em]
Analysis & 1 : 1 & AnalysisConclusion & Cada análise tem uma única conclusão formal. \[0.4em]
AccidentCase & 1 : 1 & Report & Cada processo origina um único relatório final. \\
\bottomrule
\caption{Relações e cardinalidades entre entidades}
\end{longtable}

% -------------------------------------------------
\section{Considerações de Design}

\subsection{Palavras-passe Nunca são Armazenadas em Plain Text}

A tabela \texttt{User} armazena um campo \texttt{password\_hash}. Antes de inserir ou atualizar um utilizador, a camada aplicacional deve cifrar a palavra-passe utilizando um algoritmo seguro como \textbf{bcrypt}.

\subsection{A Altura do Dano é um Valor Calculado}

A tabela \texttt{Measurement} armazena coordenadas de pixeis. A altura real em centímetros \emph{não} é armazenada na base de dados. Esta decisão foi tomada não só pelo facto de se um dia fosse necessário provar de onde veio essa medição não teria de ser feito todo o trabalho manual novamente, mas também a pensar na possibilidade da evolução do cálculo das medições. Caso a fórmula de cálculo seja refinada ou alargada com novos fatores, basta recalcular os valores a partir das coordenadas já armazenadas, sem que os dados históricos fiquem desatualizados ou inconsistentes.

\subsection{Polimorfismo de Evidências}

A entidade \texttt{Evidence} utiliza o padrão de \emph{Class Table Inheritance} (Herança por Tabela de Classe). Existe uma tabela base \texttt{Evidence} com os atributos comuns a todas as evidências, e para cada subtipo é criada uma tabela separada --- neste caso \texttt{ImageEvidence} --- que referencia a tabela base e armazena apenas os atributos específicos desse subtipo (como dimensões e metadados EXIF). Esta abordagem distingue-se do \emph{Single Table Inheritance}, onde todos os subtipos partilham uma única tabela com muitas colunas \texttt{NULL}; aqui cada tabela mantém apenas os dados que lhe são relevantes. As vantagens são: no caso de ser necessário adicionar outro tipo de evidência, o sistema continua a funcionar sem alterações à tabela base; as regras de integridade de cada subtipo podem ser aplicadas de forma rigorosa; e evita-se a proliferação de valores \texttt{NULL}.

% -------------------------------------------------
\section{Implementação em PostgreSQL}

\subsection{Exemplo de Tabelas}

\begin{lstlisting}[caption={Tabelas de Utilizador e Roles}]
CREATE TABLE Role (
  role_id     SERIAL PRIMARY KEY,
  name        VARCHAR(50)  NOT NULL UNIQUE,
  description TEXT
);

CREATE TABLE "User" (
  user_id        SERIAL PRIMARY KEY,
  name           VARCHAR(100) NOT NULL,
  email          VARCHAR(255) NOT NULL UNIQUE,
  password_hash  VARCHAR(255) NOT NULL,
  role_id        INT NOT NULL REFERENCES Role(role_id)
);

CREATE TABLE TOKENS (
  token_id    SERIAL PRIMARY KEY,
  user_id     INT NOT NULL REFERENCES "User"(user_id) ON DELETE CASCADE,
  token       VARCHAR(512) NOT NULL UNIQUE,
  created_at  TIMESTAMP NOT NULL DEFAULT NOW(),
  expires_at  TIMESTAMP NOT NULL,
  revoked     BOOLEAN NOT NULL DEFAULT FALSE
);
\end{lstlisting}

\begin{lstlisting}[caption={Tabela Central e Contexto do Acidente}]
CREATE TABLE AccidentCase (
  case_id     SERIAL PRIMARY KEY,
  user_id     INT NOT NULL REFERENCES "User"(user_id),
  created_at  TIMESTAMP NOT NULL DEFAULT NOW(),
  status      VARCHAR(50) NOT NULL DEFAULT 'open',
  description TEXT
);

CREATE TABLE WeatherConditions (
  weather_id    SERIAL PRIMARY KEY,
  case_id       INT NOT NULL UNIQUE REFERENCES AccidentCase(case_id) ON DELETE CASCADE,
  temperature   FLOAT,
  visibility    FLOAT,
  precipitation VARCHAR(50)
);

CREATE TABLE AccidentScene (
  scene_id       SERIAL PRIMARY KEY,
  case_id        INT NOT NULL UNIQUE REFERENCES AccidentCase(case_id) ON DELETE CASCADE,
  latitude       FLOAT,
  longitude      FLOAT,
  road_type      VARCHAR(50),
  road_gradient  FLOAT,
  terrain_slope  FLOAT
);
\end{lstlisting}

\begin{lstlisting}[caption={Veículos e Danos}]
CREATE TABLE Vehicle (
  vehicle_id  SERIAL PRIMARY KEY,
  case_id     INT NOT NULL REFERENCES AccidentCase(case_id) ON DELETE CASCADE,
  make        VARCHAR(100) NOT NULL,
  model       VARCHAR(100) NOT NULL,
  year        INT,
  plate       VARCHAR(20)
);

CREATE TABLE Damage (
  damage_id        SERIAL PRIMARY KEY,
  vehicle_id       INT NOT NULL REFERENCES Vehicle(vehicle_id) ON DELETE CASCADE,
  contact_zone     VARCHAR(100),
  deformation_type VARCHAR(100),
  impact_direction VARCHAR(50),
  height_from_ground FLOAT
);
\end{lstlisting}

\begin{lstlisting}[caption={Evidências -- Class Table Inheritance}]
-- Tabela base (atributos comuns a todas as evidências)
CREATE TABLE Evidence (
  evidence_id SERIAL PRIMARY KEY,
  case_id     INT NOT NULL REFERENCES AccidentCase(case_id) ON DELETE CASCADE,
  type        VARCHAR(50) NOT NULL, -- discriminador: 'image', 'video', 'document'
  uploaded_at TIMESTAMP NOT NULL DEFAULT NOW(),
  uploaded_by INT REFERENCES "User"(user_id)
);

-- Subtipo: evidência de imagem
CREATE TABLE ImageEvidence (
  evidence_id INT PRIMARY KEY REFERENCES Evidence(evidence_id) ON DELETE CASCADE,
  file_path   VARCHAR(512) NOT NULL,
  width_px    INT,
  height_px   INT,
  exif_data   TEXT
);
\end{lstlisting}

\begin{lstlisting}[caption={Análise, Medições e Conclusão}]
CREATE TABLE Analysis (
  analysis_id SERIAL PRIMARY KEY,
  case_id     INT NOT NULL REFERENCES AccidentCase(case_id) ON DELETE CASCADE,
  analyst_id  INT REFERENCES "User"(user_id),
  created_at  TIMESTAMP NOT NULL DEFAULT NOW()
);

-- Tabela de junção N:M entre Analysis e ImageEvidence
CREATE TABLE AnalysisImage (
  analysis_id INT NOT NULL REFERENCES Analysis(analysis_id) ON DELETE CASCADE,
  evidence_id INT NOT NULL REFERENCES ImageEvidence(evidence_id) ON DELETE CASCADE,
  PRIMARY KEY (analysis_id, evidence_id)
);

CREATE TABLE Measurement (
  measurement_id       SERIAL PRIMARY KEY,
  analysis_id          INT NOT NULL REFERENCES Analysis(analysis_id) ON DELETE CASCADE,
  ref_obj_x1           FLOAT, ref_obj_y1 FLOAT,
  ref_obj_x2           FLOAT, ref_obj_y2 FLOAT,
  damage_area_x1       FLOAT, damage_area_y1 FLOAT,
  damage_area_x2       FLOAT, damage_area_y2 FLOAT
);

CREATE TABLE DamageComparison (
  comparison_id SERIAL PRIMARY KEY,
  analysis_id   INT NOT NULL REFERENCES Analysis(analysis_id) ON DELETE CASCADE,
  damage_id_a   INT NOT NULL REFERENCES Damage(damage_id),
  damage_id_b   INT NOT NULL REFERENCES Damage(damage_id)
);

CREATE TABLE AnalysisConclusion (
  conclusion_id      SERIAL PRIMARY KEY,
  analysis_id        INT NOT NULL UNIQUE REFERENCES Analysis(analysis_id) ON DELETE CASCADE,
  compatibility      VARCHAR(50) NOT NULL,
  justification      TEXT
);

CREATE TABLE Report (
  report_id   SERIAL PRIMARY KEY,
  case_id     INT NOT NULL UNIQUE REFERENCES AccidentCase(case_id) ON DELETE CASCADE,
  generated_at TIMESTAMP NOT NULL DEFAULT NOW(),
  file_path   VARCHAR(512)
);
\end{lstlisting}

\vfill
\begin{center}
\small\textcolor{gray}{IDAP -- Insurance Damage Assessment Platform \quad|\quad
ISEL Projeto e Seminário 2025/26 \quad|\quad
Afonso Santos \& Rodrigo Meneses}
\end{center}

\clearpage
\appendix
\section{Diagrama EA}
\label{sec:appendix-er}

\begin{figure}[H]
    \centering
    \includegraphics[width=\textwidth]{User-Centric Accident Case-2026-03-26-175523.png}
    \caption{Modelo EA do sistema IDAP.}
\end{figure}

\vfill
\begin{center}
    \small\textcolor{gray}{[Fim da Documentação]}
\end{center}

\end{document}